\section{Introduction}

Text-to-image diffusion models, and Stable Diffusion~\cite{Rombach_2022_CVPR} in
particular, have achieved widespread deployment while relying on massive web-scraped
corpora such as Re-LAION-5B~\cite{Schuhmann_2022_Laion5b}. The automated scale of web scraping
introduces systematic biases that reflect and potentially amplify societal
inequities~\cite{Birhane_2021_MultiModalDatasetsMisogynyPornography}. Although
demographic bias in generative outputs is well documented~\cite{Luccioni_2023_AnalyzingSocietalRepresentationsDiffusionModels, Bianchi_2023_TTIGenAmplifiesDemographicStereotypes}, existing research treats models
as black boxes, measuring outcome distributions without exposing the underlying
mechanisms. A critical yet under explored dimension involves spurious features; these are low-level non-semantic image characteristics such as colour statistics, texture patterns, edge distributions, and compression artefacts that correlate
with sensitive demographic attributes and enable biased shortcuts~\cite{Meister_2023_GenderArtifactsinVisualDatasets, Zeng_2024_UnderstandingBiasinLargeScaleVisualDatasets}. Existing mitigations typically require costly retraining or fine-tuning~\cite{Friedrich_2023_FairDiffusionInstructingTextToImage}, creating barriers for practitioners working with frozen models.

This thesis addresses these gaps through a mechanistic investigation of spurious
features in Re-LAION-5B and Stable Diffusion v1.5 (SDv1.5), applying established
detection methodologies to expose low-level features that distinguish dataset sources
and correlate with demographic attributes, then evaluating targeted inference-time
interventions. Five objectives structure the work: (\textit{O1})~dataset acquisition
and demographic annotation across 200 ISCO-08-aligned occupations;
(\textit{O2})~spurious feature detection via classifier-based transformation
analysis; (\textit{O3})~demographic correlation assessment and comparison against
real-world occupational statistics; (\textit{O4})~inference-time debiasing of
SDv1.5 without model retraining; and (\textit{O5})~synthesis of findings on how
training data biases propagate through the generative process.